%%%%%%%%%%%%%%%%%%%%%%%%%%%%%%%%%%%%%%%%%%%%%%%%%%%%%%%%%%%%%%%%%%%%%%%%%%%%%%%
% COVER LETTER + JOURNAL-FIRST NOVELTY STATEMENT (IEEE TSE submission)
% Editor-facing document. Upload to the ScholarOne "Cover letter / Comments"
% slot (editor-only; not shown to reviewers, not published).
% Compile: pdflatex cover_letter_submission.tex   (single run, no bibtex)
%%%%%%%%%%%%%%%%%%%%%%%%%%%%%%%%%%%%%%%%%%%%%%%%%%%%%%%%%%%%%%%%%%%%%%%%%%%%%%%
\documentclass[11pt]{article}
\usepackage[utf8]{inputenc}
\usepackage[T1]{fontenc}
\usepackage[a4paper,margin=1in]{geometry}
\usepackage{parskip}
\usepackage{enumitem}
\usepackage{hyperref}
\hypersetup{hidelinks}
\pagestyle{plain}

\begin{document}

%============================ COVER LETTER ============================
\noindent\textbf{\large Cover Letter}\\[2pt]
\noindent IEEE Transactions on Software Engineering

\medskip
\noindent\textbf{Date:} May 20, 2026

\noindent\textbf{To:} Editor-in-Chief, IEEE Transactions on Software
Engineering, IEEE Computer Society

\noindent\textbf{From:} Acauan C. Ribeiro, Corresponding Author, on behalf of
all authors. Instituto de Computação (IComp), Universidade Federal do Amazonas
(UFAM), Manaus, AM, Brazil. \href{mailto:acauan.ribeiro@icomp.ufam.edu.br}{acauan.ribeiro@icomp.ufam.edu.br}

\hrulefill

\medskip
Dear Editor-in-Chief,

On behalf of all authors, I am pleased to submit our manuscript entitled
\textbf{``Filo-Priori: Co-Failure Graph Attention for Test Case Prioritization
in Continuous Integration''} for consideration as a \textbf{Regular (Journal
First)} manuscript in \textit{IEEE Transactions on Software Engineering}.

The manuscript addresses Test Case Prioritization (TCP) in Continuous
Integration (CI), a topic that fits TSE's interest in software testing,
software assessment methods, and empirical studies with potential impact on
software engineering practice. Existing TCP methods often score test cases as
independent entities. Filo-Priori instead models behavioral relationships
between tests using GATv2 attention over a co-failure test relationship graph,
complemented by a DNN ensemble with validation-optimized alpha blending for
metadata-sparse settings.

\noindent The main contributions are:
\begin{enumerate}[topsep=2pt,itemsep=1pt]
  \item A graph-based TCP framework centered on co-failure relationships, with
  an optional multi-edge extension when richer CI metadata is available.
  \item An empirical evaluation against five strong baselines on two
  complementary datasets: an industrial CI pipeline with 52,102 executions
  across 1,339 builds and RTPTorrent with 20 open-source Java projects and
  2,937 builds with failures.
  \item Evidence that Filo-Priori achieves APFD = 0.761 on the industrial
  dataset with significant improvements over all five baselines, and the
  highest numerical APFD of 0.854 on RTPTorrent, where gains over NodeRank and
  RETECS remain robust after Holm-Bonferroni correction.
  \item A negative result showing that generic Sentence-BERT embeddings do not
  provide a statistically significant improvement when structural features and
  co-failure propagation are available, supported by a Random-Fixed control
  experiment.
  \item Ablation, temporal validation, and sensitivity analyses showing that
  co-failure graph modeling is the decisive component and clarifying when
  graph-based TCP is most beneficial.
\end{enumerate}

This manuscript reports new research results that have not been previously
published in a peer-reviewed venue, posted as a preprint, submitted as a
conference/workshop paper, or disseminated as a technical report or thesis
chapter. It is not under consideration by any other journal, conference, or
workshop, and it will not be submitted elsewhere while under review by TSE. A
200-word Journal First Novelty Statement is included on the next page.

A public replication package, including the source code for Filo-Priori and
all baselines, trained model weights, training configurations, and scripts to
reproduce the RTPTorrent experiments end-to-end, is available at
\url{https://github.com/filotransformer/filopriori}. The RTPTorrent dataset
used for the cross-domain evaluation is publicly available under CC BY 4.0 from
its original authors. The industrial QTA dataset is proprietary and governed by
a non-disclosure agreement; it cannot be redistributed, so the industrial
experiments are not externally reproducible. This is disclosed in the
manuscript.

All authors have approved this submission. Funding, conflicts of interest, and
the use of generative AI assistance are disclosed in the manuscript. In
particular, authors Jose Carlos Rangel do Nascimento and Nícolas Riccieri
Gardin Assumpção are employed by Motorola Mobility, which provided the
industrial dataset and partially supported the research under the Brazilian
Federal Law No. 8.387/1991. Motorola Mobility had no role in the study design,
data analysis, interpretation of results, or decision to submit this manuscript
for publication.

We also disclose related prior work by overlapping authors: one of our five
baselines, FailRank-BB, is prior work co-authored by a subset of the present
authors. In this submission we treat it strictly as an external, independently
reimplemented baseline. Filo-Priori's contributions (the co-failure GATv2
framework, the validation-optimized DNN-ensemble blending, and the negative
result on generic semantic embeddings) are distinct from, and not reported in,
that prior work, and no text, figures, or results are reused from it. We note
this so that any authorship overlap surfaced during IEEE's integrity screening
is understood in context.

If ScholarOne requests suggested or non-preferred reviewers, we will provide
them in the corresponding submission fields and will avoid individuals with
actual, perceived, or potential conflicts of interest.

Thank you for considering our manuscript for review in \textit{IEEE
Transactions on Software Engineering}.

\medskip
\noindent Sincerely,

\noindent\textbf{Acauan Cardoso Ribeiro}\\
Corresponding Author, on behalf of all co-authors\\
\href{mailto:acauan.ribeiro@icomp.ufam.edu.br}{acauan.ribeiro@icomp.ufam.edu.br}\\
Instituto de Computação (IComp), Universidade Federal do Amazonas (UFAM),
Manaus, AM, Brazil

\newpage

%======================= NOVELTY STATEMENT =======================
\noindent\textbf{\large Novelty Statement: Regular (Journal First) Submission}

\medskip
\noindent\textbf{Manuscript:} Filo-Priori: Co-Failure Graph Attention for Test
Case Prioritization in Continuous Integration

\noindent\textbf{Authors:} Acauan C. Ribeiro, Eduardo L. Feitosa, Andre L. da
Costa Carvalho, Eulanda M. dos Santos, Bruno F. Gadelha, Yan R. Soares, Jose
Carlos Rangel do Nascimento, Nícolas Riccieri Gardin Assumpção

\noindent\textbf{Target Venue:} IEEE Transactions on Software Engineering
\quad\textbf{Article Type:} Regular (Journal First)

\hrulefill

\medskip
\noindent\textbf{Statement of Novelty (200 words).}
This manuscript reports completely new research results for Test Case
Prioritization (TCP) in Continuous Integration and is not a previously
published, posted, or submitted conference/workshop extension. No preliminary
version, technical report, thesis chapter, or preprint has been disseminated.
The submission is not under consideration elsewhere.

The novelty is substantive relative to IEEE TSE's Journal First criteria.
Filo-Priori introduces a TCP framework that treats tests as related entities
rather than independent instances, using GATv2 attention over a co-failure test
relationship graph and a DNN ensemble with validation-optimized alpha blending
for metadata-sparse settings. The study also contributes an explicit negative
result: generic Sentence-BERT embeddings do not significantly improve
prioritization when structural features and co-failure propagation are
available (p = 0.309), and a Random-Fixed control confirms that the apparent
semantic signal behaves as a stable identity proxy (p = 0.965). Empirically,
the work evaluates five strong baselines across an industrial CI pipeline
(52,102 executions, 1,339 builds) and RTPTorrent (20 Java projects, 2,937
failure builds). Finally, ablation, temporal validation, sensitivity analysis,
and per-edge-type experiments identify co-failure edges as the decisive graph
component and clarify the boundary conditions under which graph-based TCP
outperforms flat-feature models. These results are new, complete, and central
to the paper.

\medskip
\noindent\textbf{Confirmation of originality.} The authors confirm that the
manuscript contains new contributions to the field of Test Case Prioritization.
The work has not been published, posted as a preprint, submitted elsewhere, or
derived from a preliminary conference/workshop version. Should the work be
accepted by IEEE TSE, the authors understand that any ICSE Journal-First
presentation invitation would be handled under the journal's and conference's
applicable policies.

\medskip
\noindent\textbf{Contact (Corresponding Author):} Acauan C. Ribeiro
(\href{mailto:acauan.ribeiro@icomp.ufam.edu.br}{acauan.ribeiro@icomp.ufam.edu.br}),
Instituto de Computação (IComp), Universidade Federal do Amazonas (UFAM),
Manaus, AM, Brazil.

\end{document}
